\par\Needspace{4\baselineskip}
\begin{center}
\SourceLeaf{030b}\pdfbookmark[2]{12.}{article-012}12.
\end{center}
A suitable approximate value of $E$ from which to begin the calculation will usually be fairly apparent, especially when the problem is to be solved for several values of $M$ and has already been solved for some of them. If no other aids are available, we at least know that $E$ must lie between the limits $M$ and $M \pm e$ (with the eccentricity $e$ expressed in arcseconds and the upper sign taken in the first and second quadrants, the lower in the third and fourth). Thus, as the initial value of $E$, we may take either $M$ itself or a value obtained by increasing or decreasing it by a rough estimate. It hardly needs saying that, whenever we begin from a fairly inaccurate value, the first calculation calls for no great precision, and that smaller tables such as those prepared by the celebrated \textsc{Lalande} are amply sufficient. Moreover, to simplify the calculation we shall always choose values of $\epsilon$ whose sines can be taken directly from the tables without interpolation. These can, for example, be values in whole arcminutes or whole multiples of ten arcseconds, depending on whether the tables used proceed by single arcminutes or by ten arcseconds at a time. Anyone can work out how to adapt these directions if the angles are expressed according to the new decimal division.
